\usepackage[utf8]{inputenc}
\usepackage[german]{babel}
\usepackage{url}
\usepackage{nicefrac}
%\usepackage[export]{adjustbox} % Used for vertical alignment of tikz images, but failed
\usepackage{graphicx}
\usepackage{tikz,tikzscale}
\usepackage{pgfplots}
\pgfplotsset{compat=1.18}
\usetikzlibrary{arrows.meta,calc}
\usepackage{amsthm}
\usepackage{float}
\usepackage{environ}
\usepackage{listings}
\lstset{language=Matlab}


\title{Mathematik für Molekulare Biotechnologie B}
\author{Guido Kanschat}
\date{Fakultät für Ingenieurwissenschaften\\Universität Heidelberg}

\usepackage[mode=multiuser]{fixme}
%\fxsetup{innerlayout=inline}
%\fxsetup{targetlayout=colorcb}
\fxusetheme{color}
\FXRegisterAuthor{gk}{envgk}{GK}
\FXRegisterAuthor{mr}{envmr}{MR}

\theoremstyle{remark}
\newenvironment{Aufgabe}{\begin{onlyenv}<article>\subsubsection{Aufgabe}}{\end{onlyenv}}
\newenvironment{Wdh}{\begin{onlyenv}<article>\textbf{Wiederholung: }}%
{\end{onlyenv}}
\newenvironment{Vertiefung}{\begin{onlyenv}<article>\textbf{Vertiefung: }}%
{\end{onlyenv}}
\newenvironment{Denkpause}{%
  \begin{onlyenv}<presentation>%
    \begin{exampleblock}{Denkpause}%
}{%
    \end{exampleblock}%
  \end{onlyenv}%
}

\NewEnviron{Abschluss}{%
  \begin{frame}[label=Abschluss]
    \frametitle{Abschlussfragen \only<presentation>{\thesection.}\thesubsection}
    \begin{enumerate}
      \BODY
    \end{enumerate}%
  \end{frame}%
}

\let\epsilon\varepsilon
\let\phi\varphi
\let\theta\vartheta
\def\ii{\mathrm i}
\def\ds{\,ds}
\def\dt{\,dt}
\def\vecsym#1{\mathbf{#1}}
\def\vx{\vecsym x}
\def\vy{\vecsym y}
\def\vb{\vecsym b}
\def\vd{\vecsym d}
\def\ve{\vecsym e}
\def\vf{\vecsym f}
\def\vF{\vecsym F}
\def\vg{\vecsym g}
\def\vS{\vecsym S}
\def\vT{\vecsym T}
\def\vu{\vecsym u}
\def\vv{\vecsym v}
\def\vw{\vecsym w}
\def\R{\mathbb R}
\def\Rnn{\R^{n\times n}}

\def\norm#1{\left\|#1\right\|}
\def\scal(#1,#2){\left<#1,#2\right>}